\NativeTitle{97}{期末综合素养评价（二）}{}{90}{5}{一 & 二 & 三 & 四 & 五}{ &  &  &  & }
\NativeSlot{284}{38}{%
\ExamGroup{一、填空。{\normalfont\kaishu\fontsize{10.8bp}{14bp}\selectfont （每空1分，共19分）}}

% p097_s01

\ExamQuestion{1.}{5}%
2025年山东省的粮食播种面积是\(\underline{126339000}\)亩，横线上的数改写成用“亿”作单位的数是\mbox{（\hspace{1.8em}）}亩，保留一位小数大约是\mbox{（\hspace{1.8em}）}亿亩。\par
}
\NativeSlot{355}{46}{%
% p097_s02

\ExamQuestion{2.}{5}%
\(3.925\)由\mbox{（\hspace{1.8em}）}个\(1\)、\mbox{（\hspace{1.8em}）}个\(0.1\)、\mbox{（\hspace{1.8em}）}个\(0.01\)和\mbox{（\hspace{1.8em}）}个\(0.001\)组成，请用▲在数线上表示出这个数。\par
\BookFigure{figures/p097_s02-01.tex}
}
\NativeSlot{433}{36}{%
% p097_s03

\ExamQuestion{3.}{5}%
小军在做小数加法时，把其中一个加数\(3.5\)的小数点忽略了，算出来的结果比原来多\mbox{（\hspace{1.8em}）}。\par
}
\NativeSlot{482}{36}{%
% p097_s04

\ExamQuestion{4.}{5}%
\(32\times(\bigcirc+\triangle)=32\times\underline{\quad}+32\times\underline{\quad}\)，这是运用了\(\underline{\qquad}\)律。\par
}
\NativeSlot{506}{61}{%
% p097_s05

\ExamQuestion{5.}{5}%
下面各组直线中，相交的有\mbox{（\hspace{1.8em}）}，互相垂直的有\mbox{（\hspace{1.8em}）}，互相平行的有\mbox{（\hspace{1.8em}）}。\par
\BookFigure{figures/p097_s05-01.tex}
}
\NativeSlot{630}{95}{%
\par\noindent%
\ExamTile{.74}{%
% p097_s06

\ExamQuestion{6.}{5}%
如图，将一张长\(10\,\mathrm{cm}\)、宽\(4\,\mathrm{cm}\)的长方形和一张三角形纸交叉摆放，重叠部位是一个梯形，它的高是\mbox{（\hspace{1.8em}）}厘米。如果\(\angle1=102^\circ\)，那么\(\angle2=(\quad)^\circ\)。\par
}\hfill%
\ExamTile{.235}{%
\BookFigure{figures/p097_s06-01.tex}
}\par

\par\noindent%
\ExamTile{.74}{%
% p097_s07

\ExamQuestion{7.}{5}%
学校计划在一块周长\(10\,\mathrm m\)的等腰梯形空地上修建花坛（如右图）。现要将这块梯形空地划分成一个平行四边形和一个三角形的花坛，请你在图中划分一下。\par
}\hfill%
\ExamTile{.235}{%
\BookFigure{figures/p097_s07-01.tex}
}\par
}
